\documentclass[10.5pt]{article}

\usepackage[margin=0.8in]{geometry}
\usepackage{enumitem}
\usepackage{titlesec}
\usepackage[hidelinks]{hyperref}
\usepackage{parskip}

\titleformat{\section}{\large\bfseries}{}{0em}{}[{\vspace{-4pt}\rule{\linewidth}{0.5pt}}]
\titlespacing{\section}{0pt}{3pt}{3pt}
\setlist[itemize]{leftmargin=1.25em, itemsep=2pt, topsep=2pt}
\pagestyle{empty}
\setlength{\parskip}{4pt}

\newcommand{\pub}[1]{\textnormal{\textit{#1}}}

\begin{document}

%========================= HEADER =========================
\begin{center}
    {\LARGE\bfseries Mainul Islam}\\[4pt]
    {\large Research Summary}
\end{center}

\vspace{2pt}

%========================= INTERESTS =========================
\section*{Research Interests}
\begin{itemize}
    \item Learning-based and adaptive control for autonomous systems --- when learned policies beat
    well-tuned model-based control, and how to certify it.
    \item Robust perception, sensor fusion, and state estimation under delayed, intermittent, or
    drifting measurements.
    \item Sim-to-real transfer and digital twins --- closing the simulate--design--validate gap
    systematically.
    \item Efficient, interpretable ML for safety-critical edge deployment.
\end{itemize}

%========================= PUBLICATIONS =========================
\section*{Publications}
\begin{itemize}
    \item \pub{MonoSIM} (\textit{IEEE ICARA 2026}, B.Sc.\ thesis) --- open-source software-in-the-loop
    framework for Ackermann vehicles. \textit{I built the simulation testbed and the monocular
    sliding-window lane-detection pipeline, then implemented and benchmarked MPC against PID for path
    tracking} --- closing the perception-to-control loop entirely in simulation. The platform is released
    openly so other groups can validate control strategies before committing to hardware.
    \item \textit{Solar Energy} (Elsevier, IF 6.6, 2025) --- two-stage PV fault inspection.
    \textit{I developed the deep-learning segmentation and fault-classification pipeline (U-Net/RefineNet,
    93.6\% IoU; MobileNetV3, 97.5\% accuracy) and added LIME-based explainability} to make each diagnosis
    auditable for field maintenance. The framework generalizes across two public datasets and diverse
    real-world fault types, supporting automated, transparent PV monitoring.
    \item \textit{Energy Conversion and Management: X} (Elsevier, IF 8.8, 2026) --- topological survey of
    zeta power converters, from classical designs to emerging architectures. \textit{I contributed to the
    topology comparison and analysis.} It consolidates scattered converter designs into a single comparative
    reference to guide topology selection.
    \item \textit{IEEE ICAD 2026} --- robust Bangla license-plate recognition.
    \textit{I contributed the adaptive YOLOv8 detection stage (97.8\% accuracy) and the out-of-distribution
    evaluation on a curated low-light dataset}, coupled to a ViT\,+\,BanglaBERT vision--language OCR reader.
    The vision--language design handles conjunct characters and multi-line layouts that defeat conventional
    OCR engines.
    \item \textit{IEEE ICECIE 2024} --- long-horizon electricity-demand forecasting.
    \textit{I built and benchmarked the ML forecasters (Random Forest, KNN, XGBoost, LightGBM) and the
    feature pipeline} on a million-sample utility dataset ($R^2 = 0.94$). The comparison frames model choice
    as a practical planning tool for a fast-growing developing-country grid.
    \item \textit{IEEE ICCIT 2024} --- distributed real-time embedded networking.
    \textit{I co-developed the master--slave RS485 communication firmware and its timing analysis across
    eleven STM32 nodes}, validated on a 6-DOF manipulator ROV. A selective-feedback scheme keeps every node
    synchronized without saturating the half-duplex bus, enabling scalable actuator and sensor networks.
    \item \textit{Under review} --- speed synchronization of a dual-BLDC tracked differential-drive
    vehicle using a GA-tuned multi-rate cross-coupled controller. \textit{I co-developed the control scheme
    and its evaluation.} The approach targets accurate straight-line tracking and reduced heading drift on
    skid-steer platforms.
\end{itemize}

%========================= SYSTEMS =========================
\section*{Research-Grade Systems Built and Deployed}
As an R\&D engineer at Cybernetics Hi-Tech Solutions, I own fielded platforms end-to-end --- the source of
most of my research questions.
\begin{itemize}
    \item \textbf{Autonomous fleet (CYBERFLEET).} Complete ROS\,2 stack for a three-robot AGV fleet:
    five-state EKF fusing 50\,Hz odometry with AprilTag observations (Joseph-form updates, out-of-sequence
    handling), space-time A\textsuperscript{*} multi-agent planning, and a 50\,Hz trajectory tracker.
    \item \textbf{Man-transportable EOD robot (7-DOF).} Closed-form inverse kinematics verified against the
    URDF and integrated with MoveIt\,2; CANopen (CiA-402) control of eleven servo drives over two CAN buses,
    with a layered fail-safe architecture.
    \item \textbf{Onboard-vision UAV.} YOLO detection with DeepSORT/KCF tracking on embedded accelerator
    hardware --- a study in accuracy-versus-latency trade-offs under power constraints.
\end{itemize}

\end{document}
